\subsection{Convolution and transposition}

Let us return to the hypotheses and notations of the beginning of No.~1, but let us assume in addition that \(\beta\) is relatively invariant with multiplier \(\chi\) ; \(\chi\) is therefore a continuous function on G.

PROPOSITION 7. --- Let \(f\) be a locally \(\beta\) -integrable function on \(X\) , \(\nu\) a measure on \(X\) , and \(\mu\) a measure on \(G\) . Assume that:

\begin{enumerate}
\def\labelenumi{(\roman{enumi})}
\item
  \(\mu\) and \(f\) are convolvable and formula (3) of No.~1 defines locally \(\beta\) -almost everywhere a convolution product \(\mu *^{\beta} f\) .
\item
  \(\chi \cdot \check{\mu}\) and \(\nu\) are convolvable.
\item
  The function \(g(s,x) = f(s^{-1}x)\chi (s^{-1})\) is \((\mu \otimes \nu)\) -integrable.
\end{enumerate}

Then \(f\) is essentially integrable for \((\chi \cdot \breve{\mu}) * \nu\) , the function \(\mu *^{\beta} f\) defined by (3) is \(\nu\) -integrable, and

\[
\nu (\mu * ^ {\beta} f) = \big ((\chi \cdot \check {\mu}) * \nu \big) (f).\tag{7}
\]

Since \(g(s, x)\) is integrable for \(\mu \otimes \nu\) , the function \(f(sx)\) is essentially integrable for \((\chi \cdot \breve{\mu}) \otimes \nu\) and \(f\) is essentially integrable for \((\chi \cdot \breve{\mu}) * \nu\) . By the Lebesgue-Fubini theorem, \(\mu *^{\beta} f = \int g(s, x) d\mu(s)\) is \(\nu\) -integrable and

\[
\begin{array}{l} \nu (\mu * ^ {\beta} f) = \iint f (s ^ {- 1} x) \chi (s ^ {- 1}) d \mu (s) d \nu (x) \\ = \iint f (s x) \chi (s) d \stackrel {\vee} {\mu} (s) d \nu (x) = \big ((\chi \cdot \stackrel {\vee} {\mu}) * \nu \big) (f). \end{array}
\]

Examples. --- 1) One can take \(f \in \mathcal{C}(\mathrm{X})\) , \(\nu \in \mathcal{C}'(\mathrm{X})\) and \(\mu \in \mathcal{C}'(\mathrm{G})\) by Prop. 3, and the Cor. of Prop. 5 of §1, No.~4. The formula (7) then means that the endomorphism \(\nu \mapsto (\chi \cdot \breve{\mu}) * \nu\) of \(\mathcal{C}'(\mathrm{X})\) is the transpose of the endomorphism \(f \mapsto \mu * f\) of \(\mathcal{C}(\mathrm{X})\) .

\begin{enumerate}
\def\labelenumi{\arabic{enumi})}
\setcounter{enumi}{1}
\item
  One can take \(f \in \mathcal{K}(\mathrm{X})\) , \(\nu \in \mathcal{M}(\mathrm{X})\) and \(\mu \in \mathcal{C}'(\mathrm{G})\) by Prop. 3, Prop. 8 of §3, No.~2, and the remark that the support of the continuous function \(g(s,x)\) intersects the support of \(\mu \otimes \nu\) in a compact set. The formula (7) then means that the endomorphism \(\nu \mapsto (\chi \cdot \breve{\mu}) * \nu\) of \(\mathcal{M}(\mathrm{X})\) is the transpose of the endomorphism \(f \mapsto \mu * f\) of \(\mathcal{K}(\mathrm{X})\) .
\item
  If G operates properly on X, one can take \(f \in \mathcal{K}(X)\) , \(\nu \in \mathcal{C}'(X)\) and \(\mu \in \mathcal{M}(G)\) by Prop. 4, Prop. 8 of §3, No.~2, and the same remark as in Example 2.
\end{enumerate}

PROPOSITION 8. --- Let \(f\) and \(g\) be two locally \(\beta\) -integrable functions on \(X\) and let \(\mu \in \mathcal{M}(G)\) . Assume that:

\begin{enumerate}
\def\labelenumi{(\roman{enumi})}
\item
  \(\mu\) and \(f\) are convolvable and the formula (3) of No.~1 defines locally \(\beta\) -almost everywhere a convolution product \(\mu *^{\beta} f\) .
\item
  \(\chi \cdot \breve{\mu}\) and \(g\) are convolvable and the formula (3) of No.~1 (with \(\mu\) replaced by \(\chi \cdot \breve{\mu}\) and \(f\) by \(g\) ) defines locally \(\beta\) -almost everywhere a convolution product \((\chi \cdot \breve{\mu}) *^{\beta} g\) .
\item
  There exists a function \(\psi\) on \(\mathbf{G}\) , equal locally \(\mu\) -almost everywhere to 1, such that the function
\end{enumerate}

\[
h (s, x) = g (x) f \left(s ^ {- 1} x\right) \chi \left(s ^ {- 1}\right) \psi (s)
\]

is \((\mu \otimes \beta)\) -integrable.

Then the functions \(g(x)\big((\mu * ^ {\beta} f)(x)\big)\) and \(f(x)\big(\big((\chi \cdot \breve{\mu}) * ^ {\beta} g\big)(x)\big)\) are essentially \(\beta\) -integrable, and

\[
\int f (x) \big (\big ((\chi \cdot \breve {\mu}) * ^ {\beta} g \big) (x) \big) d \beta (x) = \int g (x) \big ((\mu * ^ {\beta} f) (x) \big) d \beta (x).\tag{8}
\]

For, by (iii) and the Lebesgue--Fubini theorem, the function

\[
x \mapsto g (x) \int f (s ^ {- 1} x) \chi (s ^ {- 1}) \psi (s) d \mu (s)
\]

is \(\beta\) -integrable, and

\[
\begin{array}{l} \mathrm{I} = \iint f (s ^ {- 1} x) g (x) \chi (s ^ {- 1}) \psi (s) d \mu (s) d \beta (x) \\ = \int g (x) d \beta (x) \int f (s ^ {- 1} x) \chi (s ^ {- 1}) \psi (s) d \mu (s). \end{array}
\]

But \(\psi \cdot \mu = \mu\) , consequently

\[
\int f (s ^ {- 1} x) \chi (s ^ {- 1}) \psi (s) d \mu (s) = (\mu * ^ {\beta} f) (x)
\]

locally \(\beta\) -almost everywhere. This shows that the function

\[
x \mapsto g (x) \left(\left(\mu * ^ {\beta} f\right) (x)\right)
\]

is essentially \(\beta\) -integrable and that

\[
\mathrm{I} = \int g (x) \left(\left(\mu * ^ {\beta} f\right) (x)\right) d \beta (x).
\]

On the other hand, Lemma 1 shows that the function

\[
(s, x) \mapsto g (s x) f (x) \chi (s ^ {- 1}) \psi (s)
\]

is integrable for \((\chi \cdot \mu) \otimes \beta\) . Therefore the function \((s, x) \mapsto g(s^{-1}x)f(x)\psi(s^{-1})\) is integrable for \(\breve{\mu} \otimes \beta\) , and

\[
\begin{array}{l} \mathrm{I} = \iint g (s ^ {- 1} x) f (x) \psi (s ^ {- 1}) d \mu^ {\vee} (s) d \beta (x) \\ = \int f (x) d \beta (x) \int g (s ^ {- 1} x) \psi (s ^ {- 1}) d \mu^ {\vee} (s). \end{array}
\]

But \(\check{\psi} \cdot \check{\mu} = \check{\mu}\) and so \(\int g(s^{-1}x)\psi(s^{-1})d\check{\mu}(s) = ((\chi \cdot \check{\mu}) *^{\beta}g)(x)\) locally \(\beta\) -almost everywhere. This shows that the function

\[
x \mapsto f (x) \left(\left((\chi \cdot \stackrel {\vee} {\mu}) * ^ {\beta} g\right) (x)\right)
\]

is essentially \(\beta\) -integrable and that

\[
\mathrm{I} = \int f (x) \left(\left(\left(\chi \cdot \stackrel {\vee} {\mu}\right) * ^ {\beta} g\right) (x)\right) d \beta (x).
\]

This proves the proposition.

Examples. --- 4) One can take \(f \in \mathcal{C}(\mathrm{X})\) , \(g \in \mathcal{K}(\mathrm{X})\) and \(\mu \in \mathcal{C}'(\mathrm{G})\) (with \(\psi = 1\) ).

\begin{enumerate}
\def\labelenumi{\arabic{enumi})}
\setcounter{enumi}{4}
\item
  If \(\mathbf{G}\) operates properly on \(\mathbf{X}\) , one can take \(f \in \mathcal{K}(\mathbf{X})\) , \(g \in \mathcal{K}(\mathbf{X})\) and \(\mu \in \mathcal{M}(\mathbf{G})\) (with \(\psi = 1\) ).
\item
  One can take \(f \in \mathrm{L}^p(\mathrm{X},\beta)\) , \(g \in \mathrm{L}^q(\mathrm{X},\beta)\) and \(\mu \in \mathcal{M}^\rho(\mathrm{G})\) , where \(1 \leqslant p < +\infty\) , \(\frac{1}{p} + \frac{1}{q} = 1\) , \(\rho = \chi^{-1/q}\) . The conditions (i) and (ii) are satisfied by Props. 5 and 6. Let us prove (iii). We have seen that \(\mu\) is carried by a set \(S\) that is a countable union of compact sets. Let us take for \(\psi\) the characteristic function of \(S\) . The function \(h\) is \((\mu \otimes \beta)\) -measurable: for, the function \((s,x) \mapsto g(x)\chi(s^{-1})\psi(s)\) is so, as is the function \((s,x) \mapsto f(s^{-1}x)\) by Lemma 1. Moreover, \(g\) being zero outside a countable union of \(\beta\) -integrable sets, \(h\) is zero outside a countable union of \((\mu \otimes \beta)\) -integrable sets. We then have (Ch. V, §8, No.~3, Prop. 7):
\end{enumerate}

\[
\begin{array}{l} \mathrm{J} = \iint^ {*} | g (x) f (s ^ {- 1} x) | \chi (s ^ {- 1}) \psi (s) d | \mu | (s) d \beta (x) \\ = \int^ {*} | g (x) | d \beta (x) \int^ {*} | f (s ^ {- 1} x) | \chi (s ^ {- 1}) \psi (s) d | \mu | (s). \end{array}\tag{9}
\]

But since \(g\) (resp. \(\psi\) ) is zero outside a countable union of integrable sets, the upper integrals of the second member of (9) are equal to the essential upper integrals (Ch. V, §1, No.~2, Prop. 7). Now (Ch. V, §5, No.~3, Prop. 3)

\[
\int^ {\bullet} | f (s ^ {- 1} x) | \chi (s ^ {- 1}) \psi (s) d | \mu | (s) = \int^ {\bullet} | f (s ^ {- 1} x) | \chi (s ^ {- 1}) d | \mu | (s)
\]

since \(\mu = \psi \cdot \mu\) . By Prop. 6, this last integral is finite and is equal to \((|\mu| *^{\beta} |f|)(x)\) locally \(\beta\) -almost everywhere. Therefore

\[
J = \int^ {\bullet} | g (x) | \left(| \mu | * ^ {\beta} | f |\right) (x) d \beta (x),
\]

and J is finite since \(g \in \mathbf{L}^q\) and \(|\mu| *^\beta |f| \in \mathbf{L}^p\) (Prop. 6). Therefore \(h\) is \((\mu \otimes \beta)\) -integrable.

The formula (8) then means that the endomorphism \(g \mapsto (\chi \cdot \breve{\mu}) * g\) of \(\mathrm{L}^q(\mathrm{X},\beta)\) is, for \(\mu \in \mathcal{M}^\rho(\mathrm{G})\) , the transpose of the endomorphism \(f \mapsto \mu * f\) of \(\mathrm{L}^p(\mathrm{X},\beta)\) .

